\section{第 14.1 节：自旋的引入与本章路线图}
\label{section:QM-C14-S01}
\studymeta
  {QM-C14-S01}
  {2026-10-27 至 2026-10-28}
  {Shankar, \textit{Principles of Quantum Mechanics}, Second Edition}
  {第 14 章 \S 14.1 Introduction；印刷 p.~373 / PDF p.~380，至同页 \S 14.2 标题前}
  {一块讲解与概念核对已完成；笔记已确认}

\notetopic{QM-C14-S01-T01}{补充新的自由度，而非放弃量子框架}
\keyidea{Spin（自旋）是没有直接经典对应物的量子自由度；仅将已有经典位置、动量变量正则量子化，并不能自动得到电子自旋。}

\textbf{引入动机。}此前常从经典物理量建立量子算符，例如由 $\mathbf l=\mathbf r\times\mathbf p$ 得到轨道角动量 $\mathbf L=\mathbf R\times\mathbf P$。这条路线以经典模型已包含相应自由度为前提。自旋没有这种直接对应，因而不能机械地照搬经典表达式来确定自旋算符及其 Hamiltonian；还需结合物理事实、已知量子结构与类比建立描述。

\dialoguesupplement{教材所说“不能直接应用四条公设”不等于公设失效。态矢量、可观测量算符和 Hamiltonian 决定演化的框架仍适用；需要扩充的是状态空间及其中的物理内容，而非把原有位置、动量公式算得更精细。}

\textbf{本章随后解决的三个问题。}
\begin{center}
\small
\renewcommand{\arraystretch}{1.15}
\begin{tabularx}{\textwidth}{@{}p{4.6cm}X@{}}
  \toprule
  层次 & 要建立的内容 \\
  \midrule
  自旋的物理性质 & 自旋是什么，以及它与已有角动量描述的关系。 \\
  Kinematics（运动学） & 如何扩充 Hilbert 空间，自旋算符如何作用；先建立状态、算符及变换规则，不是研究粒子的运动轨迹。 \\
  Dynamics（动力学） & 自旋如何进入 Hamiltonian，以及由此产生的时间演化。 \\
  \bottomrule
\end{tabularx}
\end{center}

\textbf{范围。}本章主要研究电子自旋；光子不能按这里的非相对论方式处理，教材留至第 18 章。本节仅作导读，不提前展开 \S 14.2 起的物理性质与数学构造；无独立编号 Example 或 Exercise，不另设题组。

\textbf{一分钟回顾。}自旋要求引入新的量子自由度，但不推翻量子力学的基本框架。后续主线是：明确物理性质，建立状态空间与算符，再研究动力学。

\section{第 14.2 节：自旋是内禀角动量}
\label{section:QM-C14-S02}
\studymeta
  {QM-C14-S02}
  {2026-10-28 至 2026-10-29}
  {Shankar, \textit{Principles of Quantum Mechanics}, Second Edition}
  {第 14 章 \S 14.2 What Is the Nature of Spin?；印刷 pp.~373 下部--374 上部 / PDF pp.~380--381，至 \S 14.3 标题前；含印刷 p.~374 的相关脚注}
  {一块讲解与概念核对已完成；笔记已确认}

\notetopic{QM-C14-S02-T01}{轨道运动之外的角动量}
\keyidea{Spin（自旋）是 intrinsic angular momentum（内禀角动量），不来自粒子在空间中的轨道运动，也不能用经典小球自转模型来解释。}

\textbf{为什么仅有轨道角动量不够？}已有的算符为
\begin{equation}
  \mathbf L=\mathbf R\times\mathbf P=-i\hbar\,\mathbf r\times\nabla.
\end{equation}
教材考虑理想的零动量电子，其空间波函数为常数。因此 $\nabla\psi=0$，从而 $L_x\psi=L_y\psi=L_z\psi=0$。若全部角动量都由 $\mathbf L$ 描述，沿任意轴的测量结果就只能为零；但电子即使没有轨道角动量，仍具有额外的角动量。把这部分记作 $\mathbf S$，其沿选定轴的分量有两个可能测量值：
\begin{equation}
  \boxed{S_z\text{ 的可能测量值为 }+\frac{\hbar}{2}\text{ 或 }-\frac{\hbar}{2}.}
\end{equation}
因此，轨道角动量为零不能推出全部角动量为零；需要引入新的自旋自由度，而非修改上述轨道算符的计算。

\textbf{教材脚注的限定。}实际实验通常测量与角动量相关的 magnetic moment（磁矩）；自旋最初是由光谱证据发现的，并非上述理想化实验。正文的描述用于突出旧的轨道自由度不足，不应当作发现过程的历史复述。

\textbf{保留名称，不保留机械模型。}早期曾把电子的额外角动量想象成陀螺式自转，因而使用 spin 这一名称；教材保留名称，但不采用相应的经典机械模型。研究自旋可以从它是一种角动量这一物理事实出发，借助已学的量子角动量结构建立描述，无须假设电子内部各部分正在绕轴运动。

\dialoguesupplement{常数空间波函数是广义动量本征态，不能在整个空间按普通方式归一化。$\pm\hbar/2$ 表示沿某一轴的分量，不是角动量大小有正负；仅知道两个可能结果也不能断言它们等概率，概率还取决于自旋状态。}

\textbf{一分钟回顾。}自旋是轨道运动之外的内禀角动量，不是经典自转。本节确定其物理性质；状态空间与算符的具体构造留到 \S 14.3。
